\documentclass[12pt]{article}
\usepackage[T1]{fontenc}
\usepackage[utf8]{inputenc}
\usepackage{lmodern}
\usepackage{textcomp}
\usepackage{enumitem}
\usepackage{geometry}
\geometry{margin=1in}
\begin{document}
\begin{center}
Answer Key - Religious Studies - K-12 Level Assessment 12 \
\small (Answers are ordered exactly as in the question paper.)
\end{center}

\section*{Multiple Choice}
\begin{enumerate}[label=\arabic*.]
\item \textbf{Answer:} B
\\ \textit{Options:} A) Worship; B) Submission; C) Obedience; D) Faith
\\ \textit{Note:} The word "Islam" in Arabic is derived from the root word "salam," which means peace, and it signifies submission to the will of Allah, reflecting the core principle of the religion where followers submit to God's guidance and authority.
\item \textbf{Answer:} D
\\ \textit{Options:} A) Bishop Oscar Romero; B) Resistance and renewal; C) Dignity of all human beings; D) All of the above
\\ \textit{Note:} The Communidades de Base, rooted in liberation theology, encompasses various aspects such as grassroots organization, social justice activism, and community engagement, making "All of the above" the correct answer as it reflects the multifaceted nature of these communities.
\item \textbf{Answer:} A
\\ \textit{Options:} A) Bardo Thodol; B) Bardo Ngawang; C) Bardo Gyaltsen; D) Bardo Chogye
\\ \textit{Note:} The Tibetan Book of the Dead is called "Bardo Thodol" because this term translates to "Liberation Through Hearing in the Intermediate State," which reflects its purpose of guiding individuals through the process of dying and the experiences of the afterlife in Tibetan Buddhism.
\item \textbf{Answer:} A
\\ \textit{Options:} A) Gregory; B) Augustine; C) Athanasius; D) Francis
\\ \textit{Note:} Pope Gregory the Great promoted the veneration of St. Peter and St. Paul to strengthen the papacy's authority and unify the Christian community in Rome during a time of political and social upheaval.
\item \textbf{Answer:} C
\\ \textit{Options:} A) The Holy Book; B) The Narrative; C) The Recitation; D) The Pillars
\\ \textit{Note:} The term "Qur'an" is derived from the Arabic verb "qara'a," which means "to recite," making "The Recitation" the literal translation and highlighting the text's emphasis on oral transmission and recitation in Islamic practice.
\end{enumerate}

\section*{True / False}
\begin{enumerate}[label=\arabic*.]
\setcounter{enumi}{5}
\item \textbf{Answer:} False
\\ \textit{Options:} A) True; B) False
\\ \textit{Note:} The term 'Panj Kakke' refers to the five articles of faith that Sikhs are required to wear, not the sacred texts of Sikhism, which are known as the Guru Granth Sahib.
\item \textbf{Answer:} False
\\ \textit{Options:} A) True; B) False
\\ \textit{Note:} Mahayana Buddhism is not the dominant form of Buddhism in Sri Lanka; instead, Theravada Buddhism is the predominant tradition practiced in the country.
\item \textbf{Answer:} False
\\ \textit{Options:} A) True; B) False
\\ \textit{Note:} The Mandate of Heaven was actually established during the Zhou dynasty, not the Shang dynasty, to justify the Zhou's overthrow of the Shang and to legitimize their own rule.
\item \textbf{Answer:} True
\\ \textit{Options:} A) True; B) False
\\ \textit{Note:} The Continental Reformation began in 1517 when Martin Luther famously posted his Ninety-Five Theses, challenging the practices of the Catholic Church and sparking a movement that significantly altered Christian beliefs and practices across Europe.
\item \textbf{Answer:} False
\\ \textit{Options:} A) True; B) False
\\ \textit{Note:} Sacred literature in Jainism predates Mahavira, as it was established by earlier Tirthankaras and passed down through oral tradition before being formally written.
\end{enumerate}

\section*{Long Form Response}
\begin{enumerate}[label=\arabic*.]
\setcounter{enumi}{10}
\item \textbf{Answer:} In the Egyptian composition Ludul Bel Nemequi, the character of Marduk is depicted in ways that highlight both his wrath and mercy. Marduk's wrath is evident in his readiness to punish wrongdoers and maintain order in the cosmos, reflecting the belief that chaos must be confronted with strength. Conversely, he also shows mercy by offering forgiveness and protection to those who sincerely repent and seek his guidance. This duality illustrates the complex nature of divine power, where harsh justice is balanced by compassion, ultimately emphasizing the themes of wrath and mercy in the context of maintaining harmony in the world. Such portrayals suggest that Marduk, as a deity, embodies the principles necessary for both justice and redemption.
\\ \textit{Note:} correct because it effectively captures the dual nature of Marduk's character in "Ludul Bel Nemequi," illustrating how his actions reflect the themes of wrath through punishment and chaos management, while simultaneously demonstrating mercy through forgiveness and protection for the repentant.
\item \textbf{Answer:} Jupiter is the king of the gods in Roman mythology, known for his role as the god of sky and thunder, similar to his Greek counterpart, Zeus. Both deities are depicted as powerful figures who wield lightning bolts and preside over the other gods, reflecting their positions as leaders in their respective pantheons. Jupiter's significance in Roman culture is evident in his association with law, order, and the protection of the state, often linked to the welfare of the Roman people. In contrast, Zeus also embodies these traits but is more often associated with personal relationships and human affairs in Greek mythology. Both gods symbolize authority and justice, but their cultural roles highlight the different values of Roman and Greek societies, with Jupiter emphasizing state and community while Zeus focuses on individual connections and morality.
\\ \textit{Note:} correct because it accurately identifies Jupiter as the Roman king of the gods, compares his attributes and functions to Zeus in Greek mythology, and highlights Jupiter's cultural significance in terms of law, order, and state protection, which are central to Roman values.
\end{enumerate}

\vfill
\noindent\textit{End of Answer Key}
\end{document}